\section{A meta-observation about prime ideals}
\label{section-oka-families}

\noindent
This section is taken from the CRing project. Let $R$ be a ring and
let $S \subset R$ be a multiplicative subset.
A consequence of
Lemma \ref{lemma-spec-localization}
is that an ideal $I \subset R$ maximal with respect to the property
of not intersecting $S$ is prime. The reason is that $I = R \cap \mathfrak m$
for some maximal ideal $\mathfrak m$ of the ring $S^{-1}R$.
It turns out that for many properties of ideals, the maximal ones
are prime. A general method of seeing this was developed in \cite{Lam-Reyes}.
In this section, we digress to explain this phenomenon.

\medskip\noindent
Let $R$ be a ring. If $I$ is an ideal of $R$ and $a \in R$, we
define
$$
(I : a) = \left\{ x \in R \mid xa \in I\right\}.
$$
More generally, if $J \subset R$ is an ideal, we define
$$
(I : J) = \left\{ x \in R \mid xJ \subset I\right\}.
$$

\begin{lemma}
\label{lemma-colon}
Let $R$ be a ring. For a principal ideal $J \subset R$, and for any ideal
$I \subset J$ we have $I = J (I : J)$.
\end{lemma}

\begin{proof}
Say $J = (a)$. Then $(I : J) = (I : a)$.
Since $I \subset J$ we see that any $y \in I$ is of the form
$y = xa$ for some $x \in (I : a)$. Hence $I \subset J (I : J)$.
Conversely, if $x \in (I : a)$, then $xJ = (xa) \subset I$, which
proves the other inclusion.
\end{proof}

\noindent
Let $\mathcal{F}$ be a collection of ideals of $R$. We are interested in
conditions that will guarantee that the maximal elements in the complement
of $\mathcal{F}$ are prime.

\begin{definition}
\label{definition-oka-family}
Let $R$ be a ring. Let $\mathcal{F}$ be a set of ideals of $R$. We say
$\mathcal{F}$ is an {\it Oka family} if $R \in \mathcal{F}$ and
whenever $I \subset R$ is an ideal and $(I : a), (I, a) \in \mathcal{F}$
for some $a \in R$, then $I \in \mathcal{F}$.
\end{definition}

\noindent
Let us give some examples of Oka families. The first example is the basic
example discussed in the introduction to this section.

\begin{example}
\label{example-oka-family-not-meet-multiplicative-set}
Let $R$ be a ring and let $S$ be a multiplicative subset of $R$.
We claim that $\mathcal{F} = \{I \subset R \mid I \cap S \not = \emptyset\}$
is an Oka family. Namely, suppose that $(I : a), (I, a) \in \mathcal{F}$
for some $a \in R$. Then pick $s \in (I, a) \cap S$ and
$s' \in (I : a) \cap S$. Then $ss' \in I \cap S$ and hence
$I \in \mathcal{F}$. Thus $\mathcal{F}$ is an Oka family.
\end{example}

\begin{example}
\label{example-oka-family-finitely-generated}
Let $R$ be a ring, $I \subset R$ an ideal, and $a \in R$.
If $(I : a)$ is generated by $a_1, \ldots, a_n$ and $(I, a)$
is generated by $a, b_1, \ldots, b_m$ with
$b_1, \ldots, b_m \in I$, then $I$ is generated by
$aa_1, \ldots, aa_n, b_1, \ldots, b_m$.
To see this, note that if $x \in I$, then $x \in (I, a)$
is a linear combination of $a, b_1, \ldots, b_m$, but the
coefficient of $a$ must lie in $(I:a)$.
As a result, we deduce that
the family of finitely generated ideals is an Oka family.
\end{example}

\begin{example}
\label{example-oka-family-principal}
Let us show that the family of principal ideals of a ring $R$ is an Oka family.
Indeed, suppose $I \subset R$ is an ideal, $a \in R$, and $(I, a)$ and
$(I : a)$ are principal. Note that $(I : a) = (I : (I, a))$.
Setting $J = (I, a)$, we find that $J$ is principal and $(I : J)$ is too. By
Lemma \ref{lemma-colon}
we have $I = J (I : J)$.
Thus we find in our situation that since $J = (I, a)$ and $(I : J)$
are principal, $I$ is principal.
\end{example}

\begin{example}
\label{example-oka-family-bound-cardinality}
Let $R$ be a ring.
Let $\kappa$ be an infinite cardinal.
The family of ideals which can be generated by at most $\kappa$ elements
is an Oka family. The argument is analogous to the argument in
Example \ref{example-oka-family-finitely-generated}
and is omitted.
\end{example}

\begin{example}
\label{example-oka-family-property-modules}
Let $A$ be a ring, $I \subset A$ an ideal, and $a \in A$ an element.
There is a short exact sequence $0 \to A/(I : a) \to A/I \to A/(I, a) \to 0$
where the first arrow is given by multiplication by $a$. Thus if $P$
is a property of $A$-modules that is stable under extensions and holds
for $0$, then the family of ideals $I$ such that $A/I$ has $P$
is an Oka family.
\end{example}

\begin{proposition}
\label{proposition-oka}
If $\mathcal{F}$ is an Oka family of ideals, then any maximal element of
the complement of $\mathcal{F}$ is prime.
\end{proposition}

\begin{proof}
Suppose $I \not \in \mathcal{F}$ is maximal with respect to not being in
$\mathcal{F}$ but $I$ is not prime. Note that $I \not = R$ because
$R \in \mathcal{F}$. Since $I$ is not prime we can find $a, b \in R - I$
with $ab \in I$. It follows that $(I, a) \neq I$ and $(I : a)$ contains
$b \not \in I$ so also $(I : a) \neq I$. Thus $(I : a), (I, a)$ both
strictly contain $I$, so they must belong to $\mathcal{F}$.
By the Oka condition, we have $I \in \mathcal{F}$, a contradiction.
\end{proof}

\noindent
At this point we are able to turn most of the examples above into
a lemma about prime ideals in a ring.

\begin{lemma}
\label{lemma-simple}
Let $R$ be a ring. Let $S$ be a multiplicative subset of $R$.
An ideal $I \subset R$ which is maximal with respect to the property
that $I \cap S = \emptyset$ is prime.
\end{lemma}

\begin{proof}
This is the example discussed in the introduction to this section.
For an alternative proof, combine
Example \ref{example-oka-family-not-meet-multiplicative-set}
with
Proposition \ref{proposition-oka}.
\end{proof}

\begin{lemma}
\label{lemma-cohen}
Let $R$ be a ring.
\begin{enumerate}
\item An ideal $I \subset R$ maximal with respect to not being
finitely generated is prime.
\item If every prime ideal of $R$ is
finitely generated, then
every ideal of $R$ is finitely generated\footnote{Later we will say
that $R$ is Noetherian.}.
\end{enumerate}
\end{lemma}

\begin{proof}
The first assertion is an immediate consequence of
Example \ref{example-oka-family-finitely-generated} and
Proposition \ref{proposition-oka}. For the second,
suppose that there exists an ideal $I \subset R$ which is not finitely
generated. The union of a totally ordered chain $\left\{I_\alpha\right\}$
of ideals that are not finitely generated is not finitely generated;
indeed, if $I = \bigcup I_\alpha$ were generated by
$a_1, \ldots, a_n$, then all the generators would belong to some
$I_\alpha $ and would consequently generate it.
By Zorn's lemma, there is an ideal maximal with respect to being not finitely
generated. By the first part this ideal is prime.
\end{proof}

\begin{lemma}
\label{lemma-primes-principal}
Let $R$ be a ring.
\begin{enumerate}
\item An ideal $I \subset R$ maximal with respect to not being
principal is prime.
\item If every prime ideal of $R$ is principal, then
every ideal of $R$ is principal.
\end{enumerate}
\end{lemma}

\begin{proof}
The first part follows from
Example \ref{example-oka-family-principal} and
Proposition \ref{proposition-oka}.
For the second, suppose that there exists an ideal $I \subset R$
which is not principal. The union of a totally ordered chain
$\left\{I_\alpha\right\}$ of ideals that not principal is not principal;
indeed, if $I = \bigcup I_\alpha$ were generated by
$a$, then $a$ would belong to some $I_\alpha $ and $a$ would generate it.
By Zorn's lemma, there is an ideal maximal with respect to not being
principal. This ideal is necessarily prime by the first part.
\end{proof}

\begin{lemma}
\label{lemma-characterize-domain}
Let $R$ be a ring.
\begin{enumerate}
\item An ideal maximal among the ideals which do not contain a
nonzerodivisor is prime.
\item If $R$ is nonzero and every nonzero prime ideal in $R$
contains a nonzerodivisor, then $R$ is a domain.
\end{enumerate}
\end{lemma}

\begin{proof}
Consider the set $S$ of nonzerodivisors. It is a multiplicative
subset of $R$. Hence any ideal maximal with respect to not intersecting
$S$ is prime, see
Lemma \ref{lemma-simple}.
Thus, if every nonzero prime ideal contains a nonzerodivisor, then
$(0)$ is prime, i.e., $R$ is a domain.
\end{proof}

\begin{remark}
\label{remark-cohen-bound-cardinality}
Let $R$ be a ring. Let $\kappa$ be an infinite cardinal.
By applying
Example \ref{example-oka-family-bound-cardinality} and
Proposition \ref{proposition-oka}
we see that any ideal maximal with respect to the property of not being
generated by $\kappa$ elements is prime. This result is not so
useful because there exists a ring for which every prime ideal
of $R$ can be generated by $\aleph_0$ elements, but some
ideal cannot. Namely, let $k$ be a field, let $T$ be a set whose
cardinality is greater than $\aleph_0$ and let
$$
R = k[\{x_n\}_{n \geq 1}, \{z_{t, n}\}_{t \in T, n \geq 0}]/
(x_n^2, z_{t, n}^2, x_n z_{t, n} - z_{t, n - 1})
$$
This is a local ring with unique prime ideal
$\mathfrak m = (x_n)$. But the ideal $(z_{t, n})$ cannot
be generated by countably many elements.
\end{remark}

\begin{example}
\label{example-noetherian-topology-on-spec}
\begin{reference}
Comment by Lukas Heger of November 12, 2020.
\end{reference}
Let $R$ be a ring and $X = \Spec(R)$. Since closed subsets of $X$ correspond to
radical ideas of $R$ (Lemma \ref{lemma-Zariski-topology}) we see that $X$ is a
Noetherian topological space if and only if we have ACC for radical ideals.
This holds if and only if every radical ideal is the radical of a finitely
generated ideal (details omitted). Let
$$
\mathcal{F} = \{I \subset R \mid
\sqrt{I} = \sqrt{(f_1, \ldots, f_n)}\text{ for some }n
\text{ and }f_1, \ldots, f_n \in R\}.
$$
The reader can show that $\mathcal{F}$ is an Oka family by using the identity
$$
\sqrt{I} = \sqrt{(I, a)(I : a)}
$$
which holds for any ideal $I \subset R$ and any element $a \in R$.
On the other hand, if we have a totally ordered chain of ideals
$\{I_\alpha\}$ none of which are in $\mathcal{F}$, then the union
$I = \bigcup I_\alpha$ cannot be in $\mathcal{F}$ either. Otherwise
$\sqrt{I} = \sqrt{(f_1, \ldots, f_n)}$, then $f_i^e \in I$ for some $e$,
then $f_i^e \in I_\alpha$ for some $\alpha$ independent of $i$, then
$\sqrt{I_\alpha} = \sqrt{(f_1, \ldots, f_n)}$, contradiction.
Thus if the set of ideals not in $\mathcal{F}$ is nonempty, then it
has maximal elements and
exactly as in Lemma \ref{lemma-cohen} we conclude that $X$ is a
Noetherian topological space if and only if every prime ideal of $R$
is equal to $\sqrt{(f_1, \ldots, f_n)}$ for some $f_1, \ldots, f_n \in R$.
If we ever need this result we will carefully state and prove this
result here.
\end{example}








